% =====================================================================
%  Guide de la bibliothèque timeseries
%
%  Compilation :  make guide-fr        (depuis la racine du dépôt)
%  ou, à la main :  latexmk -xelatex guide-fr.tex
%
%  xelatex est requis pour les caractères semi-graphiques du schéma du
%  float64, que seule une police système comme Menlo possède.
% =====================================================================
\documentclass[11pt,a4paper]{report}

\usepackage{fontspec}
\setmonofont{Menlo}[Scale=0.85]
\usepackage[french]{babel}
% Le texte porte déjà ses espaces avant « : » ; sans cela, babel en
% ajoute une dans « 00:57 » et coupe les heures en deux.
\frenchsetup{AutoSpacePunctuation=false}
\usepackage{csquotes}
\usepackage{microtype}

\usepackage[a4paper,margin=2.5cm]{geometry}
\usepackage{booktabs}
\usepackage{longtable}
\usepackage{array}
\usepackage{fancyvrb}
\usepackage{fvextra}
\usepackage{etoolbox}
\usepackage{calc}          % largeurs de colonnes calculées

% Macros que pandoc suppose fournies par son gabarit.
\providecommand{\tightlist}{\setlength{\itemsep}{0pt}\setlength{\parskip}{0pt}}
\providecommand{\passthrough}[1]{#1}
\providecommand{\real}[1]{#1}

% Exposants absents de Latin Modern.
\usepackage{newunicodechar}
\newunicodechar{⁵}{\textsuperscript{5}}
\newunicodechar{²}{\textsuperscript{2}}
\newunicodechar{ᵉ}{\textsuperscript{e}}
\newunicodechar{ʳ}{\textsuperscript{r}}
\newunicodechar{·}{\textperiodcentered}
\newunicodechar{±}{$\pm$}
\newunicodechar{×}{$\times$}
\newunicodechar{−}{$-$}
% Espaces de largeur nulle semées par la conversion des tableaux.
\catcode`\^^^^200b=9\relax
\usepackage{titlesec}
\usepackage{fancyhdr}
\usepackage[table,dvipsnames]{xcolor}
\usepackage{hyperref}
\hypersetup{colorlinks=true, linkcolor=black, citecolor=black,
            urlcolor=MidnightBlue,
            pdftitle={Guide de la bibliothèque timeseries},
            pdfauthor={Frédéric Flament}}

% --- Bibliographie ------------------------------------------------------
\usepackage[backend=biber, style=numeric, sorting=none]{biblatex}
\addbibresource{references.bib}

% --- Titres : alignés à gauche, sans page blanche entre les chapitres ---
\titleformat{\chapter}[hang]
  {\LARGE\bfseries}{\thechapter.}{0.6em}{}
\titlespacing*{\chapter}{0pt}{0pt}{1.2em}
\titleformat{\section}{\Large\bfseries\raggedright}{\thesection}{0.8em}{}
\titleformat{\subsection}{\large\bfseries\raggedright}{\thesubsection}{0.8em}{}
\setlength{\parindent}{0pt}
\setlength{\parskip}{0.55em}

% --- Tableaux comparatifs, souvent larges -------------------------------
\AtBeginEnvironment{longtable}{\small}
\AtBeginEnvironment{tabular}{\small}

% --- Blocs de code ------------------------------------------------------
\fvset{fontsize=\small, breaklines=true, breakanywhere=true}
\DefineVerbatimEnvironment{verbatim}{Verbatim}{fontsize=\small, breaklines=true}

% --- En-têtes -----------------------------------------------------------
\pagestyle{fancy}
\fancyhf{}
\fancyhead[L]{\small\nouppercase{\leftmark}}
\fancyhead[R]{\small\textit{timeseries}}
\fancyfoot[C]{\small\thepage}
\renewcommand{\headrulewidth}{0.4pt}
\renewcommand{\chaptermark}[1]{\markboth{\thechapter.\ #1}{}}

% report compose l'abstract sur une page à lui, précédé d'un \titlepage,
% et le coiffe du mot « Résumé ». Redéfini ici pour qu'il tienne sous le
% titre, sans intitulé : à cette place, le texte se présente seul.
% L'italique n'est pas la convention de LaTeX, qui compose l'abstract en
% corps droit réduit ; c'est celle de bien des revues.
\renewenvironment{abstract}{\begin{quotation}\small\itshape\noindent}{\end{quotation}}

% Les métadonnées du document, définies une fois. La page de titre et
% hyperref les reprennent : il n'y a pas d'autre endroit où les écrire.
\newcommand{\montitre}{Guide de la bibliothèque \textit{timeseries}}
\newcommand{\monsoustitre}{Séries temporelles irrégulières,  traitement de masse sous  \textit{high availability}, séries à faible contenu informatif}
\newcommand{\monauteur}{Frédéric Flament}
\newcommand{\monorganisation}{Développé pour Usefulrisk ---
\href{https://usefulrisk.com}{usefulrisk.com}}
% Deux régimes distincts : le code s'emprunte, le guide se lit. La
% licence MIT du dépôt porte sur la bibliothèque, pas sur ce document.
\newcommand{\malicencecode}{Bibliothèque sous licence MIT,
\textcopyright{} 2025--2026 Frédéric Flament}
\newcommand{\malicenceguide}{Ce guide : \textcopyright{} 2026
Frédéric Flament, tous droits réservés}

\begin{document}

% --- Première page : titre, auteur, date et abstract ensemble ----------
% report place le titre et l'abstract sur deux pages séparées ; ici tout
% tient sur la première, la table des matières commençant en dessous.
\begin{titlepage}

\begin{center}
  \vspace*{5cm}

  {\LARGE\bfseries \montitre}\\[0.6em]
  {\large \monsoustitre}

  \vspace{3cm}

  {\large \monauteur}\\[1em]
  {\today}
\end{center}

\vspace{3cm}

\begin{abstract}
La bibliothèque \texttt{timeseries} traite des séries de mesures matérielles, caractérisées par des intervalles irréguliers et des trous là où la chaîne de
transmission est rompue. Elle nettoie, régularise et compacte --- sans
que des relevés absents ou incomplets ne cassent la série.

Les trous dans les relevés sont traités par NaN-boxing : l'absence (Not-a-Value) est un flottant ordinaire. Une
série traverse sans convention supplémentaire tout code qui accepte
un \texttt{[]float64}, alors que les solutions usuelles --- structure à
champ booléen, pointeur nullable, valeur sentinelle ---
imposent une extraction avant chaque calcul, ou se contournent au risque
de produire des erreurs silencieuses. Le procédé est éprouvé ailleurs :
Prometheus distingue de la même façon, par la charge utile d'un NaN, une
série arrêtée d'une erreur de calcul (\S\ref{sec:tsdb}).

Les mesures tiennent dans un tableau de flottants sans pointeur, les agrégats se lisent en une passe, et le traitement par enregistrement reste
constant tant que la série tient en mémoire vive. Des mesures de performances et les
variantes écartées figurent au chapitre~\ref{chap:performances}.

La bibliothèque réduit sans perte  les séries à faible contenu informatif --- relevés à
variabilité faible ---, avec redéploiement.

Écrite en Go pour le rapport entre simplicité d'écriture/lecture et la performance.

\end{abstract}

\vfill

% Le pied de page de titre : à qui appartient la bibliothèque, et sous
% quelles conditions elle s'utilise. \vfill ci-dessus le pousse en bas
% de la page, quelle que soit la longueur de l'abstract.
\begin{center}
  \small
  \monorganisation\\[0.5em]
  \malicencecode\\[0.2em]
  \malicenceguide
\end{center}

\end{titlepage}

\tableofcontents

% Les tables portent un titre et une étiquette, donc elles se listent et
% se citent par leur numéro (\ref{tab:...}).
\listoftables

\chapter{Introduction}\label{chap:introduction}

\texttt{timeseries} est une bibliothèque Go qui traite des séries de
relevés mesurés. Elle répond à quatre questions.

\textbf{Que faire d'un relevé qui n'existe pas ?} Un trou n'est ni un
zéro, ni une erreur, ni une ligne à supprimer. La bibliothèque le
représente par un \textbf{NaV}, un NaN porteur d'un repère, que les
calculs ignorent au lieu de s'y arrêter --- là où un NaN ordinaire,
réservé aux erreurs de calcul, continue de se propager. Une moyenne
mensuelle survit à un jour manquant ; elle ne survit pas à une division
par zéro, et c'est la différence qu'il fallait pouvoir exprimer.

\textbf{Que faire d'un instrument qui n'est pas un métronome ?} Les
relevés annoncés « toutes les heures » arrivent à 00:57, 02:03, 03:00.
La régularisation les repose sur une grille de pas fixe, avec une
tolérance pour les enregistreurs qui dérivent, et des trous là où aucun
relevé n'est arrivé. Deux séries traitées au même pas deviennent alors
comparables point à point, condition de tout calcul qui les met en
relation.

\textbf{Que faire d'un signal qui ne dit presque rien ?} Une porte, une
consigne, un état, un compteur : la plupart des relevés n'apportent
aucune information, puisque la valeur n'a pas bougé. La bibliothèque les
réduit à leurs changements --- le premier relevé, le dernier, et ce qui
s'est passé entre les deux --- et sait les reconstituer à la demande.
Une bande morte permet de n'enregistrer qu'au-delà d'un écart donné,
avec une perte bornée et connue. Et comme un silence prolongé se
confond, après réduction, avec une valeur qui se maintient,
\texttt{MarkSilences} le marque comme un trou avant que l'information ne
disparaisse.

\textbf{Que faire quand il y en a beaucoup ?} Un capteur ne pose pas de
problème ; mille capteurs interrogés sur un an en posent. Les agrégats
parcourent donc leurs données une seule fois et n'allouent rien ; le
chargement d'un lot trie une fois, en \(n\log n\), là où l'insertion
point par point est quadratique --- deux ordres de grandeur d'écart sur
un million de relevés en désordre ; et un tampon réutilisable permet de
boucler sur des centaines de séries sans qu'une allocation ait lieu.
Chacun de ces choix est mesuré, les variantes écartées sont conservées
dans le dépôt avec leurs chiffres, et les mesures sont relatives à la
machine d'essai : ce sont les rapports qui se transposent, pas les
durées.

Le reste du guide détaille ces quatre réponses, les conventions de temps
qui les sous-tendent, et ce que chaque opération demande en mémoire, en
allocations et en parcours.


\chapter{NaV pour « Not a Value » : stratégie pour les données
manquantes}\label{chap:nav}

\section{Introduction}
Comment savoir qu'une donnée est manquante ? Ou plutôt, comment
caractériser, dans les calculs, dans les rapports, et dans le
déclenchement d'action, une donnée qui n'est pas là ? Une solution
courante est de déclencher une action à partir d'une absence de donnée
pendant une durée fixée à l'avance. Si on examine les solutions
apportées dans les logiciels courants, nous avons opté pour une solution
utilisant le \emph{NaN-boxing}.

\section{Définition de NaV}

La norme IEEE 754~\cite{ieee754}, qui régit les nombres flottants, laisse libre le
contenu de la mantisse d'un NaN : n'importe quelle valeur non nulle en
fait un NaN valide. Nous utilisons cette liberté pour introduire un
nombre particulier, le \textbf{NaV}, un NaN silencieux dont le bit 48 de
la mantisse est allumé. Il reste un \texttt{float64} ordinaire : un
\texttt{[]float64} n'a besoin d'aucun masque parallèle, et
\texttt{math.IsNaN} le reconnaît toujours, si bien que le code déjà
écrit pour se prémunir des NaN se prémunit aussi des NaV.

Dans tout ce guide, un \textbf{trou} désigne un relevé manquant,
c'est-à-dire un NaV : un instant où la mesure aurait dû avoir lieu et
n'a pas eu lieu.

\begin{longtable}[]{@{}
  >{\raggedright\arraybackslash}p{(\linewidth - 6\tabcolsep) * \real{0.2500}}
  >{\raggedright\arraybackslash}p{(\linewidth - 6\tabcolsep) * \real{0.2500}}
  >{\raggedright\arraybackslash}p{(\linewidth - 6\tabcolsep) * \real{0.2500}}
  >{\raggedright\arraybackslash}p{(\linewidth - 6\tabcolsep) * \real{0.2500}}@{}}
\caption{Origine d'un non-nombre et traitement qui en découle}\label{tab:origines}\\
\toprule\noalign{}
\begin{minipage}[b]{\linewidth}\raggedright
\end{minipage} & \begin{minipage}[b]{\linewidth}\raggedright
Signification
\end{minipage} & \begin{minipage}[b]{\linewidth}\raggedright
Exemple
\end{minipage} & \begin{minipage}[b]{\linewidth}\raggedright
Traitement
\end{minipage} \\
\midrule\noalign{}
\endhead
\bottomrule\noalign{}
\endlastfoot
\textbf{NaV} & Rien n'a été mesuré & Capteur hors service, ligne
absente, point rejeté & Ignoré \\
\textbf{NaN} & Un calcul a cassé & \texttt{0/0} en amont,
\texttt{log(-1)} & Se propage \\
\end{longtable}

\section{Distinction entre NaV et NaN}

\begin{itemize}
\tightlist
\item
  Le NaN a été créé pour les erreurs de calcul, une donnée manquante
  n'est pas une erreur de calcul
\item
  Une donnée manquante ne doit pas nécessairement se propager dans tous
  les calculs subséquents, or les NaN propagent l'indication de l'erreur
  de calcul
\end{itemize}

\section{Comportement attendu d'un trou}

\begin{enumerate}
\def\labelenumi{\arabic{enumi}.}
\tightlist
\item
  \textbf{Un trou n'arrête jamais un calcul.} Il est ignoré par les
  agrégats et neutre dans l'addition.
\item
  \textbf{Une erreur se propage toujours.} Quand un trou et une erreur
  se rencontrent, l'erreur l'emporte.
\item
  \textbf{Quand il ne reste rien à calculer, le résultat est un trou}
  --- jamais zéro.
\end{enumerate}

La troisième règle appelle un point d'attention. La moyenne d'une série
vide vaut NaV, pas 0. Une moyenne nulle est une affirmation sur les
données : elle dit que les relevés se compensent, ou, plus simplement,
que la température en Celsius était au point de gel. Sur une série vide,
il n'y a rien à affirmer, et répondre zéro serait une invention qu'aucun
lecteur ultérieur ne peut détecter.

\section{Règles de comportement}\label{sec:regles}

Les conséquences ne sont pas nécessairement triviales:

\begin{longtable}[]{@{}
  >{\raggedright\arraybackslash}p{(\linewidth - 4\tabcolsep) * \real{0.3333}}
  >{\raggedright\arraybackslash}p{(\linewidth - 4\tabcolsep) * \real{0.3333}}
  >{\raggedright\arraybackslash}p{(\linewidth - 4\tabcolsep) * \real{0.3333}}@{}}
\caption{Règles de comportement, cas par cas}\label{tab:regles}\\
\toprule\noalign{}
\begin{minipage}[b]{\linewidth}\raggedright
Situation
\end{minipage} & \begin{minipage}[b]{\linewidth}\raggedright
Résultat
\end{minipage} & \begin{minipage}[b]{\linewidth}\raggedright
Pourquoi
\end{minipage} \\
\midrule\noalign{}
\endhead
\bottomrule\noalign{}
\endlastfoot
Moyenne de \texttt{[1,\ 2,\ NaV,\ 3]} & \texttt{2} & Le trou est
ignoré ; le diviseur vaut 3, pas 4 \\
Moyenne de \texttt{[1,\ 2,\ NaN,\ 3]} & \texttt{NaN} & Une erreur en
amont doit rester visible \\
Moyenne de \texttt{[NaV,\ NaV]} & \texttt{NaV} & Rien à dire \\
\texttt{Sub(5,\ NaV)} & \texttt{NaV} & Une différence exige ses deux
termes ; rendre 5 reviendrait à lire le trou comme un zéro \\
\texttt{Add(NaV,\ 5)} & \texttt{5} & Une somme tolère un terme absent \\
\texttt{Mul(NaV,\ 0)} & \texttt{NaV} & Une quantité inconnue de quelque
chose reste inconnue \\
Variation après un trou & \texttt{NaV} & Un relevé de 14 après un trou
n'est pas une hausse de 14 \\
Point rejeté comme aberrant & \texttt{NaV} & Il a été mesuré mais n'est
pas cru : absent, pas cassé \\
Interpoler un \texttt{NaN} & laissé tel quel & Recouvrir une erreur d'un
nombre plausible est la façon dont un bug cesse d'être remarqué \\
\end{longtable}

L'asymétrie entre \texttt{Add} et \texttt{Sub} peut provoquer un
haussement de sourcils. La logique est la suivante: une somme accumule
des termes indépendants, et un terme absent laisse les autres intacts
--- c'est ce qui permet à une moyenne mensuelle de survivre à un jour
manquant. Une différence compare deux valeurs précises ; si l'une est
inconnue, la différence l'est aussi, et toute autre réponse fabrique une
variation que personne n'a mesurée.

Ces règles ne sont pas une invention locale. SQL les applique depuis
toujours : \texttt{sum} « computes the sum of the non-null input
values », et « sum of no rows returns null, not zero as one might
expect »~\cite{pgaggregates}. Un jour manquant n'empêche pas la somme du
mois, et un mois sans aucun relevé ne vaut pas zéro. C'est mot pour mot
la première et la troisième règle. §\ref{sec:bases} détaille ce que les
bases de données font de l'absence, et le seul point où cette
bibliothèque s'en écarte.

\section{NaN-boxing en détail}

Le NaN-boxing consiste à ranger de l'information \textbf{à l'intérieur}
d'un nombre flottant, dans les bits que la norme IEEE 754 laisse libres.
Pour comprendre où ils sont, il faut ouvrir un \texttt{float64}.

\subsection{Structure d'un flottant}

Un \texttt{float64} occupe 64 bits, répartis par la norme IEEE-754 en
trois champs :

\begin{verbatim}
 ┌─┬───────────┬────────────────────────────────────────────────────┐
 │S│  exposant │                     mantisse                       │
 └─┴───────────┴────────────────────────────────────────────────────┘
  1     11 bits                      52 bits
\end{verbatim}

La valeur ordinaire se lit « mantisse × 2\^{}exposant », avec le signe
devant. Mais la norme réserve une configuration : \textbf{quand les 11
bits d'exposant valent tous 1}, ce n'est plus un nombre.

\begin{itemize}
\tightlist
\item
  Si la mantisse vaut zéro, c'est l'infini, positif ou négatif selon le
  signe.
\item
  Si la mantisse vaut autre chose que zéro, c'est un \textbf{NaN}.
\end{itemize}

La norme ne dit pas \emph{quelle} valeur la mantisse doit prendre.
N'importe laquelle, pourvu qu'elle ne soit pas nulle, fait un NaN
parfaitement valide. Il y a 2⁵² − 1 mantisses possibles, deux signes,
soit \textbf{environ neuf millions de milliards de configurations qui
sont toutes des NaN} --- dont la moitié, celles dites silencieuses, sont
utilisables sans risque. Toutes indiscernables pour l'arithmétique, et
toutes perdues si personne ne s'en sert.

\subsection{Silencieux et signalant}\label{sec:quiet}

La norme distingue encore deux familles, par le bit de poids fort de la
mantisse :

\begin{itemize}
\tightlist
\item
  \textbf{NaN silencieux} (bit à 1) : il traverse les calculs sans
  bruit. C'est celui que rend \texttt{math.NaN()}, et celui que produit
  toute opération invalide sur un processeur courant.
\item
  \textbf{NaN signalant} (bit à 0) : il est censé déclencher une
  exception matérielle. En pratique, Go ne l'utilise pas, et la plupart
  des environnements le convertissent en silencieux dès la première
  opération. On ne s'en sert pas ici.
\end{itemize}

Il reste donc \textbf{51 bits libres} dans la mantisse d'un NaN
silencieux. C'est de la place perdue, que rien n'utilise.

\subsection{Mise en œuvre dans la bibliothèque}

\texttt{NaV} est un NaN silencieux dont \textbf{un} de ces bits libres
est allumé :

\begin{verbatim}
math.NaN()  : 0 11111111111 1000000000000000000000000000000000000000000000000000
NaV         : 0 11111111111 1000000001000000000000000000000000000000000000000000
                            ↑        ↑
                            │        └── le repère NaV (bit 48)
                            └── le bit « silencieux »
\end{verbatim}

Le test tient en une ligne : c'est un NaN, \textbf{et} le bit de repère
est allumé.

\begin{verbatim}
func IsNaV(x float64) bool {
    return math.IsNaN(x) && math.Float64bits(x)&navTag != 0
}
\end{verbatim}

Quatre conséquences en découlent, et ce sont elles qui justifient la
technique :

\begin{enumerate}
\def\labelenumi{\arabic{enumi}.}
\tightlist
\item
  \textbf{Aucun surcoût mémoire.} Un NaV est un \texttt{float64}. Un
  tableau de mesures reste un \texttt{[]float64}, sans tableau
  compagnon.
\item
  \textbf{Le code existant continue de fonctionner.}
  \texttt{math.IsNaN(NaV)} est vrai, donc tout code déjà écrit pour se
  prémunir des NaN se prémunit aussi des NaV. Rien à recompiler, rien à
  auditer.
\item
  \textbf{La distinction survit au transport binaire.} Copier,
  sérialiser en binaire, passer par \texttt{math.Float64bits} : le
  repère voyage avec la valeur, puisqu'il \emph{est} la valeur.
\item
  \textbf{Il reste 50 bits libres.} On pourrait un jour distinguer «
  jamais mesuré », « rejeté comme aberrant », « interpolé », sans rien
  casser de l'existant.
\end{enumerate}

\subsection{Contreparties}

\begin{itemize}
\tightlist
\item
  \textbf{L'arithmétique doit passer par les fonctions du paquet.} Ce
  que devient le repère à travers un \texttt{-} ordinaire dépend du
  processeur (\S\ref{sec:avertissement}). C'est la contrainte principale.
\item
  \textbf{La sérialisation texte perd le repère.} JSON n'a pas de NaN :
  tout non-nombre devient \texttt{null}, et la distinction ne subsiste
  que dans les compteurs (\S\ref{sec:json}).
\item
  \textbf{C'est une technique peu connue.} Un relecteur qui découvre la
  bibliothèque doit d'abord comprendre ce qu'il lit --- ce chapitre
  existe pour ça.
\item
  \textbf{Cela ne vaut que pour les flottants.} Le procédé tient aux bits
  libres du NaN d'IEEE 754~\cite{ieee754} ; un entier, un booléen, une
  chaîne n'ont pas d'encodage inutilisé à détourner. La bibliothèque le
  constate dès qu'elle quitte le flottant : une durée manquante est une
  valeur sentinelle, \texttt{NaDuration}, et non un repère logé dans des
  bits libres (chapitre~\ref{chap:temps}).
\end{itemize}

Marquer l'absence dans n'importe quel type demande donc un marqueur en
dehors de la valeur, sous l'une de deux formes : un champ booléen accolé
à chaque valeur, dont la section suivante détaille le prix, ou un bit par
valeur dans un masque séparé. C'est la voie des bases de données, et
§\ref{sec:bases} dit comment elles la suivent.

\section{Rejet de la structure à champ booléen}\label{sec:struct}

C'est la solution la plus répandue :

\begin{verbatim}
type Valued struct {
    Value float64
    Valid bool
}
\end{verbatim}

C'est ce que font \texttt{sql.NullFloat64} en Go,
\texttt{Option\textless{}f64\textgreater{}} en Rust, \texttt{double?} en
C\#, \texttt{Optional\textless{}Double\textgreater{}} en Java. Six
raisons de ne pas l'avoir reprise.

\subsection{Doublement de la mémoire}

\begin{verbatim}
float64                       :  8 octets
struct{ float64; bool }       : 16 octets
\end{verbatim}

Le booléen n'occupe qu'un octet, mais l'alignement en impose huit : sept
octets sont perdus par point. Sur un million de mesures, \textbf{8 Mo
deviennent 16 Mo} ; sur dix millions, 80 deviennent 160.

Un masque parallèle \texttt{[]bool} fait mieux --- 9 Mo --- mais au
prix d'un second tableau à transporter, à découper et à trier en même
temps que le premier, et qu'on oublie à la première refonte.

\subsection{Incompatibilité avec l'écosystème}

Tout ce qui calcule sur des flottants, en Go, attend un
\texttt{[]float64}: gonum, les transformées de Fourier, les
bibliothèques de statistiques, et \texttt{notavalue} lui-même.

Avec une structure, chaque appel exige d'abord une extraction --- une
boucle et une allocation proportionnelles à la série :

\begin{longtable}[]{@{}
  >{\raggedright\arraybackslash}p{(\linewidth - 4\tabcolsep) * \real{0.3333}}
  >{\raggedright\arraybackslash}p{(\linewidth - 4\tabcolsep) * \real{0.3333}}
  >{\raggedright\arraybackslash}p{(\linewidth - 4\tabcolsep) * \real{0.3333}}@{}}
\caption{Moyenne sur un million de points : extraction contre passage direct}\label{tab:extraction}\\
\toprule\noalign{}
\begin{minipage}[b]{\linewidth}\raggedright
Moyenne sur un million de points
\end{minipage} & \begin{minipage}[b]{\linewidth}\raggedright
Temps
\end{minipage} & \begin{minipage}[b]{\linewidth}\raggedright
Allocation
\end{minipage} \\
\midrule\noalign{}
\endhead
\bottomrule\noalign{}
\endlastfoot
Structure, puis extraction en \texttt{[]float64} & 1 368 µs &
\textbf{8 Mo par appel} \\
NaN-boxing, tableau passé tel quel & 623 µs & \textbf{aucune} \\
\end{longtable}

Deux fois plus lent, et 8 Mo alloués à chaque appel. Sur un traitement
qui enchaîne moyenne, médiane, écart-type et percentiles, la facture se
paie quatre fois. Les durées absolues valent pour la machine d'essai
décrite au chapitre~\ref{chap:performances} ; c'est le rapport entre les
deux lignes qui compte.

\subsection{Traitement répété à chaque calcul}

Chaque fonction qui travaille sur les mesures doit connaître la
structure et tester le booléen. Le traitement des absences se répète
partout au lieu d'être porté par la valeur elle-même.

\subsection{Contournement possible}

Le compilateur n'oblige jamais à lire \texttt{Valid}. Rien n'empêche
d'écrire \texttt{v.Value} sur un relevé absent : on obtient zéro, un
zéro qui a l'air d'une mesure et qui traverse tout le calcul sans
laisser de trace. L'oubli n'est signalé ni à la compilation ni à
l'exécution, et il se glisse dans la première fonction écrite un jour de
fatigue.

Un trou, lui, se défend tout seul. Il n'existe aucune façon de lire « la
valeur derrière le NaV » : la valeur \emph{est} le NaV, et toute
opération qui l'ignore le propage ou l'écarte selon les règles du
paquet. Le mauvais usage n'est pas rendu difficile, il est rendu
impossible.

\subsection{Pointeur nullable et ramasse-miettes}

\texttt{*float64}, avec \texttt{nil} pour l'absence, semble élégant.
Mais chaque point devient un pointeur de 8 octets \textbf{plus} la
valeur pointée quelque part ailleurs, et surtout : un tableau d'un
million de pointeurs est parcouru par le ramasse-miettes à chaque cycle,
alors qu'un \texttt{[]float64} ne contient aucun pointeur et lui
reste totalement invisible.

\subsection{Piège de la valeur sentinelle}

Coder l'absence par −999, ou par 0, est la solution la plus ancienne.
Elle fonctionne jusqu'au jour où une vraie mesure vaut −999 --- et ce
jour arrive. Le NaV, lui, ne peut pas être confondu avec une mesure :
aucune opération sur des nombres réels ne le produit.

\subsection{Tableau comparatif}

\begin{longtable}[]{@{}
  >{\raggedright\arraybackslash}p{(\linewidth - 8\tabcolsep) * \real{0.2000}}
  >{\raggedright\arraybackslash}p{(\linewidth - 8\tabcolsep) * \real{0.2000}}
  >{\raggedright\arraybackslash}p{(\linewidth - 8\tabcolsep) * \real{0.2000}}
  >{\raggedright\arraybackslash}p{(\linewidth - 8\tabcolsep) * \real{0.2000}}
  >{\raggedright\arraybackslash}p{(\linewidth - 8\tabcolsep) * \real{0.2000}}@{}}
\caption{Représentations de l'absence dans une série de flottants}\label{tab:representations}\\
\toprule\noalign{}
\begin{minipage}[b]{\linewidth}\raggedright
Solution
\end{minipage} & \begin{minipage}[b]{\linewidth}\raggedright
Mémoire par point
\end{minipage} & \begin{minipage}[b]{\linewidth}\raggedright
Compatible \texttt{[]float64}
\end{minipage} & \begin{minipage}[b]{\linewidth}\raggedright
Invisible au GC
\end{minipage} & \begin{minipage}[b]{\linewidth}\raggedright
Confusion possible
\end{minipage} \\
\midrule\noalign{}
\endhead
\bottomrule\noalign{}
\endlastfoot
\textbf{NaN-boxing} & 8 o & oui & oui & non \\
\texttt{struct\{float64;\ bool\}} & 16 o & non & oui & non \\
\texttt{[]float64} + \texttt{[]bool} & 9 o & oui, mais le masque
suit à part & oui & non \\
\texttt{*float64} & 8 o + la valeur & non & \textbf{non} & non \\
Sentinelle −999 & 8 o & oui & oui & \textbf{oui} \\
\end{longtable}

La vitesse de parcours, elle, ne départage pas : les quatre premières
tournent entre 0,6 et 0,8 ms par million de points, et l'écart est
dominé par le reste du calcul. \textbf{Ce qui décide, c'est la mémoire
et la compatibilité}, pas les nanosecondes.

\section{Le traitement de l'absence dans quelques bases de données courantes}\label{sec:bases}

La question est ancienne et elle a été tranchée ailleurs, à grande
échelle. Un moteur de base de données doit accepter l'absence dans
n'importe quel type, pas seulement dans un flottant, et il doit répondre
à deux questions distinctes : ce qu'un calcul fait d'une valeur absente,
et où le marqueur se loge. Les réponses au premier point sont
remarquablement proches de celles de cette bibliothèque ; celles du
second lui sont inaccessibles.

\subsection{Sémantique : deux régimes}

SQL sépare les opérateurs scalaires des agrégats, et les traite en sens
inverse.

Un opérateur scalaire propage : \texttt{NULL\ +\ 5} vaut NULL, et toute
comparaison avec NULL vaut \emph{unknown} plutôt que vrai ou faux ---
d'où la logique à trois valeurs et le \texttt{IS\ NULL}, qui existe
parce que \texttt{=\ NULL} ne peut rien affirmer.

Un agrégat saute : \texttt{sum} additionne les valeurs non nulles,
\texttt{avg} divise par leur nombre et non par celui des lignes,
\texttt{count(colonne)} compte les valeurs renseignées là où
\texttt{count(*)} compte les lignes. Et une somme sans aucune ligne
rend NULL, pas zéro~\cite{pgaggregates}. C'est la troisième règle de ce
chapitre, à l'identique : quand il n'y a rien à calculer, le résultat
est l'absence.

Reste \texttt{coalesce}, que SQL fournit pour qui veut le zéro, et qui
porte le même nom que la fonction correspondante de
\texttt{notavalue}.

\subsection{Point de divergence}

\texttt{Add(NaV,\ 5)} rend 5, là où \texttt{NULL\ +\ 5} vaut NULL. La
divergence est apparente : \texttt{Add} n'est pas le \texttt{+} de SQL,
c'est sa primitive d'accumulation, et l'accumulateur de SQL est
\texttt{sum}, qui saute les absents exactement de la même façon. Le
\texttt{+} scalaire de SQL a pour équivalent \texttt{Sub},
\texttt{Mul}, \texttt{Div}, qui propagent tous l'absence.

Une chose en revanche n'a pas d'équivalent : la deuxième règle. Un
calcul rompu n'y produit pas une valeur qui se propage, il lève une
erreur et la transaction s'arrête --- \texttt{0/0} est une
\emph{division by zero}, non un NaN silencieux. SQL peut se le
permettre : il a une transaction à annuler. Une bibliothèque appelée au
milieu d'un traitement de masse n'a pas ce recours, et le NaN qui se
propage est ce qui tient lieu d'erreur transportable.

\subsection{Stockage du marqueur}

\textbf{PostgreSQL} range le marqueur hors de la valeur, et non dans un
booléen accolé à elle. Chaque ligne porte, juste après son en-tête fixe,
une \emph{null bitmap} d'un bit par colonne --- 1 pour renseigné, 0 pour
NULL --- et cette carte n'est présente que si l'en-tête signale au moins
un NULL~\cite{pgstorage}. Une colonne NULL n'occupe alors aucun octet
pour sa valeur : le bit remplace la donnée au lieu de s'y ajouter.

\textbf{SQLite} n'a pas de carte, parce qu'il n'a pas de types de
colonne fixes : l'en-tête de chaque enregistrement énumère le type de
chaque valeur, et l'absence est l'un de ces types. Le type sérialisé 0
signifie « la valeur est NULL » et occupe zéro octet de
contenu~\cite{sqliteformat}. Le marqueur ne coûte donc rien de plus que
ce que toute valeur paie déjà, et fait l'économie du contenu.

\textbf{Apache Arrow}, le format en mémoire derrière une bonne part des
outils analytiques actuels, tient un tampon séparé : un bit par
emplacement, 1 pour renseigné, 0 pour absent. Avec une réserve que la
spécification énonce elle-même --- « a null value may occupy a
non-empty memory space in the data buffer », et le contenu de cet
emplacement est alors indéfini~\cite{arrowformat}. Sur un type de
largeur fixe, un absent paie donc le bit \emph{et} la place.

\begin{table}[htbp]
\centering
\small
\caption{Emplacement du marqueur d'absence selon le
système}\label{tab:marqueurs}
\begin{tabular}{@{}p{4.2cm}p{4.6cm}p{5.4cm}@{}}
\toprule
\textbf{Système} & \textbf{Où est le marqueur} & \textbf{Ce que coûte une
valeur absente} \\
\midrule
NaN-boxing & dans la valeur & rien : le flottant est le marqueur \\
PostgreSQL & 1 bit par colonne, en-tête de ligne & 1 bit, et aucun octet
de contenu \\
SQLite & type sérialisé par valeur, en-tête d'enregistrement & rien de
plus qu'une valeur présente, et aucun contenu \\
Apache Arrow & 1 bit par valeur, tampon séparé & 1 bit \emph{et}
l'emplacement, réservé \\
\texttt{struct\{float64;\ bool\}} & accolé à la valeur & 8 octets, par
alignement \\
\bottomrule
\end{tabular}
\end{table}

Deux des trois font mieux que tout ce que la table~\ref{tab:representations}
avait examiné.

\subsection{Bases de séries temporelles}\label{sec:tsdb}

Ce sont les plus proches parentes du sujet : elles ne stockent que des
couples instant-mesure, et rencontrent donc le problème de ce guide dans
sa forme pure. Aucune des deux n'a de NULL au sens relationnel. Un trou
y est l'absence de relevé, pas un relevé marqué.

\textbf{Prometheus} a pourtant eu besoin de dire autre chose que « pas
de relevé » : quand une cible disparaît, il faut distinguer « la série
s'est arrêtée » de « rien n'est encore arrivé ». Sans quoi la dernière
valeur connue continue d'être rendue pendant la fenêtre de rattrapage,
et un graphique montre une machine éteinte comme si elle fonctionnait
encore. La série est donc marquée périmée, et « if a query is evaluated
at a sampling timestamp after a time series is marked as stale, then no
value is returned for that time series »~\cite{promstaleness}.

Ce qui nous intéresse est la façon dont ce marqueur est représenté. Elle
est celle de ce chapitre~\cite{promvalue} :

\begin{verbatim}
// NormalNaN is a quiet NaN. This is also math.NaN().
NormalNaN uint64 = 0x7ff8000000000001

// StaleNaN is a signaling NaN, due to the MSB of the
// mantissa being 0.
StaleNaN uint64 = 0x7ff0000000000002

func IsStaleNaN(v float64) bool {
    return math.Float64bits(v) == StaleNaN
}
\end{verbatim}

Deux NaN distingués par leur charge utile, l'un pour l'erreur de calcul,
l'autre pour l'absence, et un prédicat pour les séparer. C'est
exactement \texttt{IsNaV} et \texttt{IsStdNaN}, atteint séparément et
déployé à très grande échelle. Le procédé décrit ici n'est donc pas une
curiosité locale.

Deux différences méritent d'être relevées. Prometheus prend un NaN
\emph{signalant} --- bit de poids fort de la mantisse à zéro, cf.
§\ref{sec:quiet} --- là où \texttt{notavalue} prend un NaN silencieux
marqué. Il peut se le permettre parce que son marqueur n'entre jamais
dans un calcul : il est filtré à la lecture. Un NaV, lui, traverse
\texttt{Add}, \texttt{Sub} et les agrégats, et doit donc être silencieux.
Et Prometheus compare la totalité des bits, ce qui ferme la porte à
d'autres charges utiles ; le masque sur un seul bit la laisse ouverte.

\textbf{InfluxDB} ne marque rien dans le stockage : un champ qui n'a pas
été écrit n'existe pas à cet instant. En revanche il tranche à la
lecture, et la clause \texttt{fill()} d'InfluxQL énumère les mêmes
choix que la table~\ref{tab:interpolation} --- \texttt{null},
\texttt{previous}, \texttt{linear}, \texttt{none}, ou une valeur
numérique~\cite{influxfill}. Deux d'entre eux se lisent comme un
avertissement : \texttt{fill(none)} supprime l'horodatage avec la valeur,
donc le trou lui-même, ce que \texttt{Regularize} refuse de faire ; et
\texttt{fill(0)} est précisément le geste contre lequel cette
bibliothèque existe.

Le point remarquable est ailleurs : le comportement par défaut
d'InfluxDB est \texttt{fill(null)}. Sur une fenêtre sans données, il
rend un horodatage et une absence, pas un zéro. La même décision que
celle de ce chapitre, prise par les deux moteurs de séries temporelles
les plus répandus.

\subsection{Passage du moteur à Go}

La représentation économique du moteur ne survit pas au transport. Un
pilote qui rend une colonne \texttt{double\ precision} nullable ne peut
pas la déposer dans un \texttt{float64} : la bibliothèque standard
refuse, et le message ne laisse pas de doute ---
\texttt{converting\ NULL\ to\ float64\ is\ unsupported}. Il faut élargir
le type, et \texttt{database/sql} le fait exactement sous la forme que
§\ref{sec:struct} écarte~\cite{gosql} :

\begin{verbatim}
type NullFloat64 struct {
    Float64 float64
    Valid   bool
}
\end{verbatim}

La forme générique ajoutée depuis, \texttt{sql.Null[T]}, a la même
structure : une valeur, un booléen. Le bit de la \emph{null bitmap}
devient donc un champ de huit octets dès la frontière, alignement
compris, et le tableau compact du moteur arrive en Go sous forme de
tranche de structures --- avec les six conséquences déjà énumérées.

C'est là que se mesure l'intérêt du procédé. Aucune des trois
représentations de la table~\ref{tab:marqueurs} n'est à la portée d'une
bibliothèque : un
moteur maîtrise sa représentation du premier octet sur disque au dernier
en mémoire, et peut donc décider que l'absence vit à côté de la donnée.
\texttt{timeseries} rend des \texttt{[]float64} à du code qu'elle ne
connaît pas. C'est cette tranche-là, sans carte à transporter à côté
d'elle, qui est la contrainte --- et le NaN-boxing est ce qui permet de
la tenir, au point précis où \texttt{database/sql} doit renoncer.

\section{Avertissement}\label{sec:avertissement}

Il ne faut pas utiliser les opérateurs ordinaires sur des valeurs susceptibles
d'être absentes. La charge utile que porte un NaN résultat est laissée
au processeur, et les processeurs ne s'accordent pas : sur certains,
\texttt{NaV\ -\ 5} ressort encore marqué comme NaV, sur d'autres le
repère est perdu. Passez par les \texttt{Add}, \texttt{Sub},
\texttt{Mul}, \texttt{Div} et les agrégats de \texttt{notavalue}, qui
font du résultat une propriété du code et non de la machine qui
l'exécute.

\chapter{Temps}\label{chap:temps}

\section{Stockage d'un instant}

Les horodatages sont des \href{https://pkg.go.dev/time\#Time}{\texttt{time.Time}},
tels que définis par la bibliothèque standard de Go~\cite{gotime}. Celle-ci a été
déterminante dans le choix du langage pour le traitement des données
temporelles.\texttt{time.Time} est un type bien développé en Go ---
fuseaux, formatage, comparaison sont dans la bibliothèque
standard\footnote{Voir §\ref{sec:unix} pour la lecture d'horloge
monotone, et la différence entre \emph{monotonic clock} et \emph{wall
clock}.}. Le cadre est l'utilisation de Unix Time dans toute la chaîne de traitement, le formatage en RFC3339 appliqué partout là où le passage par le texte est nécessaire.

\section{Temps Unix, sous les horodatages}\label{sec:unix}

Sous la surface, tous les systèmes qui échangent des dates s'accordent
sur une même origine : \textbf{le 1ᵉʳ janvier 1970 à 00:00:00 UTC}. Un
instant s'y ramène à un seul nombre, le compte de ce qui s'est écoulé
depuis --- secondes, millisecondes ou nanosecondes selon la précision
retenue. C'est le \emph{temps Unix}, et c'est ce que \texttt{time.Time}
manipule en interne, ce que PostgreSQL stocke, ce que transportent les
capteurs.

Quelques repères pour lire un tel nombre :

\begin{longtable}[]{@{}lll@{}}
\caption{Lectures du temps Unix}\label{tab:unix}\\
\toprule\noalign{}
Nombre & Unité & Instant \\
\midrule\noalign{}
\endhead
\bottomrule\noalign{}
\endlastfoot
\texttt{0} & seconde & 1ᵉʳ janvier 1970, 00:00:00 UTC \\
\texttt{1\ 000\ 000\ 000} & secondes & 9 septembre 2001 \\
\texttt{1\ 767\ 225\ 600} & secondes & 1ᵉʳ janvier 2026 \\
\texttt{1\ 767\ 225\ 600\ 000\ 000\ 000} & nanosecondes & le même
instant \\
\end{longtable}

Trois propriétés méritent d'être connues, parce qu'elles expliquent des
choix de la bibliothèque.

\textbf{Le temps Unix ne connaît pas les fuseaux.} C'est un compte
depuis une origine, donc un instant absolu. Deux capteurs, l'un à
Luxembourg et l'autre à Tokyo, qui mesurent au même moment produisent le
même nombre. Le fuseau n'intervient qu'à l'affichage --- et c'est
exactement pourquoi une grille régulière s'aligne dessus (\S\ref{sec:utc}) : c'est
la seule référence que deux séries partagent, où qu'elles aient été
enregistrées.

\textbf{Il ignore les secondes intercalaires.} La rotation de la Terre
n'est pas régulière, et l'UTC y ajoute de temps en temps une seconde ---
la dernière en 2016. Le temps Unix, lui, fait comme si elles
n'existaient pas : une journée y compte toujours exactement 86 400
secondes. Les systèmes qui doivent rester à l'heure les absorbent en
étirant imperceptiblement leur horloge sur quelques heures. Conséquence
pratique : une durée calculée entre deux instants séparés par une
seconde intercalaire est fausse d'une seconde, ce qui n'a d'importance
que pour la métrologie fine.

\textbf{Il ne dit rien de la précision de la mesure.} Un horodatage à la
nanoseconde n'implique pas que le capteur sache ce qu'il faisait à la
nanoseconde près. Le nombre est exact, la mesure ne l'est pas forcément
--- et la bibliothèque conserve ce qu'on lui donne sans prétendre
l'améliorer.

En Go, un \texttt{time.Time} porte davantage qu'un simple compte :
l'instant absolu, un fuseau, et parfois une
\href{https://pkg.go.dev/time\#hdr-Monotonic_Clocks}{lecture d'horloge
monotone} --- insensible aux changements d'heure du système, là où la
lecture d'horloge murale les subit --- que la bibliothèque standard
utilise pour mesurer des durées. Les horodatages
venus d'une base ou d'un capteur n'en ont pas ; c'est sans conséquence
ici.

\section{RFC 3339, forme écrite}\label{sec:rfc}

Le temps Unix est un nombre ; il faut aussi une forme écrite, lisible
par un humain et non ambiguë pour une machine. C'est la \textbf{RFC
3339}, et c'est elle que produit \texttt{ToJSON} :

\begin{verbatim}
2026-01-15T14:30:00Z           ← en UTC, le « Z » pour zéro décalage
2026-01-15T15:30:00+01:00      ← le même instant, vu de Paris
2026-01-15T14:30:00.123456789Z ← avec ses nanosecondes
\end{verbatim}

La forme est stricte, et c'est ce qui en fait la valeur : une date, un
\texttt{T}, une heure, un décalage. Du plus grand au plus petit,
toujours, avec des zéros de remplissage. C'est un sous-ensemble
volontairement réduit de la norme ISO 8601, qui autorise elle une foule
de variantes --- les semaines, les durées, les dates partielles,
l'omission des séparateurs --- et qu'aucune implémentation ne couvre
entièrement.

Trois propriétés justifient de s'y tenir :

\textbf{Le tri lexicographique est le tri chronologique.} Deux
horodatages RFC 3339 exprimés dans le même décalage se comparent comme
du texte ordinaire, caractère par caractère, et l'ordre obtenu est le
bon. C'est ce qui permet de trier un fichier de journaux avec
\texttt{sort}, ou d'indexer une colonne texte sans la convertir.

\textbf{Le décalage est obligatoire.} Un horodatage sans décalage ---
comme en produisent tant de bases de données et d'API --- ne désigne pas
un instant : \texttt{2026-01-15\ 14:30:00} peut être quatorze heures et
demie à Paris, à Tokyo ou à New York, soit trois instants distants de
plusieurs heures. La RFC 3339~\cite{rfc3339} l'interdit. Quand une donnée arrive dans
cette forme, il faut lui adjoindre le fuseau que le fournisseur
sous-entend, et c'est une décision, pas une conversion.

\textbf{Mais un décalage n'est pas un fuseau.} \texttt{+01:00} dit de
combien l'heure locale s'écarte de l'UTC à cet instant précis ; il ne
dit pas qu'on est à Paris, ni si l'heure d'été s'appliquait. Une série
sérialisée en RFC 3339 puis relue perd donc le nom de son fuseau : elle
garde l'instant exact, ce qui suffit à tout calcul, mais un regroupement
calendaire effectué après ce trajet retombera sur un décalage figé
plutôt que sur un vrai fuseau. \textbf{Quand les journées locales
comptent, regroupez avant de sérialiser}, pas après.

\section{Durée qui n'existe pas}

Le premier point d'une série n'a pas de prédécesseur, donc l'intervalle
qui le précède n'existe pas. C'est \texttt{NaDuration}, le pendant de
NaV pour le temps. Il s'affiche « NaDuration » plutôt que sous la forme
absurde de −2562047h47m16s, et \texttt{IsNaDuration} le reconnaît.

\section{Deux conventions différentes}

\begin{longtable}[]{@{}
  >{\raggedright\arraybackslash}p{(\linewidth - 6\tabcolsep) * \real{0.2500}}
  >{\raggedright\arraybackslash}p{(\linewidth - 6\tabcolsep) * \real{0.2500}}
  >{\raggedright\arraybackslash}p{(\linewidth - 6\tabcolsep) * \real{0.2500}}
  >{\raggedright\arraybackslash}p{(\linewidth - 6\tabcolsep) * \real{0.2500}}@{}}
\caption{Conventions de fenêtre : régularisation et regroupement}\label{tab:conventions}\\
\toprule\noalign{}
\begin{minipage}[b]{\linewidth}\raggedright
\end{minipage} & \begin{minipage}[b]{\linewidth}\raggedright
Fenêtre
\end{minipage} & \begin{minipage}[b]{\linewidth}\raggedright
Un relevé sur la frontière
\end{minipage} & \begin{minipage}[b]{\linewidth}\raggedright
L'instant émis
\end{minipage} \\
\midrule\noalign{}
\endhead
\bottomrule\noalign{}
\endlastfoot
\texttt{Regularize} & fermée à droite & appartient à la fenêtre qui s'y
\textbf{termine} & le top \\
\texttt{Downscale*} & ouverte à gauche & \textbf{ouvre} la nouvelle
journée & le dernier instant de la période \\
\end{longtable}

\texttt{Regularize} convient au calcul : toutes ses fenêtres durent
exactement la même chose, donc deux points pèsent toujours pareil.
\texttt{Downscale} convient au calendrier : un relevé pris à minuit
appartient à la journée qui commence, comme le dirait n'importe qui.

Les deux datent leur point à la \textbf{fin} de la période, si bien que
tous deux se lisent « tout ce qui précède jusqu'ici ».

\section{Alignement d'une grille régulière sur l'UTC}\label{sec:utc}

La grille ne s'aligne pas sur le premier relevé --- ce qui rendrait deux
séries incomparables --- mais sur le temps absolu, c'est-à-dire sur
l'UTC.

Dans un fuseau décalé d'un nombre entier d'heures, c'est invisible :

\begin{verbatim}
Paris (UTC+01:00), relevé à 10:47 → fenêtre horaire à 10:00 locales
\end{verbatim}

Dans un fuseau décalé d'une demi-heure, ça ne l'est pas :

\begin{verbatim}
Calcutta (UTC+05:30), relevé à 10:47 → fenêtre horaire à 10:30 locales
\end{verbatim}

Ce sont les mêmes instants en UTC, et c'est précisément l'objectif :
deux séries régularisées au même pas tombent sur les mêmes instants, où
qu'elles aient été enregistrées. Si ce sont les frontières d'heures
locales qui comptent --- un rapport quotidien pour une équipe sur place
---, utilisez \texttt{DownscaleDaily}, qui travaille dans le calendrier,
ou décalez les horodatages avant de régulariser.

\section{Périodes calendaires: durées irrégulières}

Les frontières du calendrier n'ont de sens qu'en un lieu : la famille
\texttt{Downscale} les calcule donc dans le fuseau du premier relevé de
la série.

Elle traite les deux jours de l'année où les horloges changent --- 23
heures au printemps, 25 à l'automne --- et les fuseaux où \textbf{minuit
n'existe pas} : à Santiago, La Havane et aux Açores, les horloges
avancent à minuit, et la journée s'ouvre à 01:00. Bâtir une période sur
un minuit qui n'a jamais eu lieu la fermerait une heure trop tôt, et
toutes les périodes suivantes avec elle.

\section{Précision}

Les statistiques sur les horodatages sont calculées sur des écarts au
premier point, et non sur des nanosecondes depuis 1970.

Un \texttt{float64} porte 53 bits de mantisse ; un horodatage Unix
actuel en nanosecondes en demande 61. Calculer sur des valeurs absolues
arrondit à quelques centaines de nanosecondes --- de quoi faire tomber
au mauvais endroit l'instant moyen de trois relevés espacés de deux
nanosecondes. Compter depuis le début de la série garde les nombres
petits et le résultat exact.

\chapter{Série}\label{chap:serie}

\section{Types}

\begin{verbatim}
type Datum struct {                 // un relevé
    Chron time.Time
    Meas  float64
}

type DataUnit struct {              // un relevé placé dans une série
    Datum
    Dchron time.Duration            // temps écoulé depuis le point précédent
    Dmeas  float64                  // variation depuis le point précédent
}
\end{verbatim}

Une \texttt{TimeSeries} est une suite de \texttt{DataUnit}, plus un
\texttt{ID}, un \texttt{Name} pour les humains et un \texttt{Comment}
qui retrace sa provenance.

\section{Invariant}

\textbf{À tout instant, les points sont dans l'ordre chronologique et
chaque delta s'accorde avec cet ordre.}

Aucune méthode ne trie une série ni ne recalcule ses deltas : il n'y a
jamais rien à réparer. \texttt{Add} et \texttt{AddBatchData}
maintiennent l'invariant en insérant --- y compris quand un relevé
arrive en retard, ce qui décale la fin du tableau et recalcule
exactement deux deltas.

D'où le champ privé. Livrer le tableau permettrait d'ajouter hors ordre
ou de trier par mesure, laissant les deltas décrire un ordre disparu. La
version précédente portait pour cela un drapeau, \texttt{deltasValid},
qui prévenait que son propre état pouvait mentir. Rendre l'état
impossible vaut mieux que le signaler.

La lecture passe par \texttt{Len}, \texttt{At}, \texttt{First},
\texttt{Last}, \texttt{Range}, \texttt{Meas} et \texttt{MeasTo}.

\section{Chargement et son coût}\label{sec:chargement}

\begin{longtable}[]{@{}ll@{}}
\caption{Chargement d'un million de relevés en désordre}\label{tab:chargement}\\
\toprule\noalign{}
Un million de relevés en désordre & Temps \\
\midrule\noalign{}
\endhead
\bottomrule\noalign{}
\endlastfoot
\texttt{Add}, un par un & \textasciitilde2,5 minutes \\
\texttt{AddBatchData} & \textasciitilde1 seconde \\
\end{longtable}

\texttt{Add} décale la fin du tableau pour chaque relevé qui doit
s'insérer plus tôt, ce qui est quadratique. Mesuré : 61 ms pour 20 000
points, 255 ms pour 40 000, 1 014 ms pour 80 000 --- un quadruplement à
chaque doublement.

\texttt{AddBatchData} ajoute tout, trie une fois et remplit les deltas
en une passe : 9, 21 et 45 ms sur les mêmes lots. Quand le lot prolonge
déjà la série dans l'ordre --- une requête avec \texttt{ORDER\ BY} ---,
il le constate en une passe et saute le tri.

\texttt{Add} pour le relevé qui arrive seul d'un flux,
\texttt{AddBatchData} pour tout ce qu'on a déjà en main.

Ces durées ont été relevées sur la machine d'essai du
chapitre~\ref{chap:performances}. Ce qui se transporte d'une machine à
l'autre n'est pas la seconde ni les deux minutes et demie, mais le
quadruplement à chaque doublement d'un côté et le doublement de l'autre.

\section{Alimenter les agrégats sans allouer}

\begin{verbatim}
var buf []float64
for _, ts := range all {
    buf = ts.MeasTo(buf[:0])           // aucune allocation après la première série
    fmt.Println(ts.Name, nav.Mean(buf))
}
\end{verbatim}

\texttt{Meas()} alloue un tableau neuf et reste la forme commode ;
\texttt{MeasTo} remplit un tampon que vous gardez, ce qui compte dès
qu'on boucle sur des centaines de séries.

\chapter{Statistiques de base}\label{chap:resume-stats}

\texttt{Stats()} renvoie trente champs, calculés à la demande et jamais
mis en cache --- un résumé ne peut donc jamais décrire une série qui a
changé depuis, et deux appels s'accordent toujours.

Les noms suivent une clé : \texttt{Ch*} concernent les horodatages,
\texttt{Ms*} les mesures, \texttt{DCh*} et \texttt{DMs*} les deltas, et
\texttt{ChAt*} donnent l'instant auquel une autre statistique se
produit.

Regroupés par la question à laquelle ils répondent :

\begin{longtable}[]{@{}
  >{\raggedright\arraybackslash}p{(\linewidth - 2\tabcolsep) * \real{0.5000}}
  >{\raggedright\arraybackslash}p{(\linewidth - 2\tabcolsep) * \real{0.5000}}@{}}
\caption{Champs de \texttt{BasicStats}, par question posée}\label{tab:basicstats}\\
\toprule\noalign{}
\begin{minipage}[b]{\linewidth}\raggedright
Question
\end{minipage} & \begin{minipage}[b]{\linewidth}\raggedright
Champs
\end{minipage} \\
\midrule\noalign{}
\endhead
\bottomrule\noalign{}
\endlastfoot
Sur quoi la série s'étend-elle ? & \texttt{Chmin}, \texttt{Chmax}, avec
\texttt{ValAtChmin} et \texttt{ValAtChmax} \\
Où commencent les données ? & \texttt{ChFirstUsable},
\texttt{ValAtFirstUsable} \\
Où se situent les points ? & \texttt{Chmean}, \texttt{Chmed} \\
Qu'a-t-on mesuré ? & \texttt{Msmin}, \texttt{Msmax} avec leurs instants,
\texttt{Msmean}, \texttt{Msmed}, \texttt{Msstd} \\
À quelle cadence arrivent les relevés ? & \texttt{DChmin},
\texttt{DChmax} avec leurs instants, \texttt{DChmean}, \texttt{DChmed},
\texttt{DChstd} \\
De combien ça bouge ? & \texttt{DMsmin}, \texttt{DMsmax},
\texttt{DMsmean}, \texttt{DMsmed}, \texttt{DMsstd} \\
Quelle est la qualité des données ? & \texttt{Len}, \texttt{NbreOfNaN},
\texttt{NbreOfNaV} \\
\end{longtable}

\textbf{\texttt{Chmin} face à \texttt{ChFirstUsable}.} Le premier dit où
la fenêtre s'ouvre, le second où les relevés commencent. Sur un capteur
qui chauffait encore, les deux diffèrent, et les confondre date la série
trop tôt.

\textbf{\texttt{NbreOfNaN} moins \texttt{NbreOfNaV}} donne le nombre de
véritables erreurs de calcul. Quand les deux sont égaux, tout ce qui
n'est pas un nombre dans la série est un trou, et rien n'est cassé.

\textbf{\texttt{DChstd}} mesure la régularité de l'échantillonnage :
proche de zéro sur un flux discipliné, il grimpe dès qu'il bégaie. C'est
souvent le nombre le plus instructif du tableau, parce qu'un intervalle
moyen masque complètement une dérive --- dix-neuf intervalles de 56 à 66
minutes ont une moyenne d'exactement une heure.

Une statistique qui ne peut pas être calculée vaut NaV, jamais zéro. Un
relevé unique n'a pas d'écart-type ; une série de trous n'a pas de
moyenne.

\chapter{Nettoyage}\label{chap:nettoyage}

\section{Définition d'un rejet}

Un point rejeté par une fonction de nettoyage n'est ni une erreur ni une suppression et produit une série temporelle dédiée.
Parfois c'est dans les données rejetées que se trouve l'information: fréquence des pannes, maintenance préventive, déclenchement d'alarmes.


Les relevés rejetés reviennent dans une série à part, avec leurs valeurs
d'origine.

\section{Valeurs jamais rejetées}

\begin{itemize}
\tightlist
\item
  \textbf{Un trou} : il n'y avait rien à juger.
\item
  \textbf{Une valeur cassée} : le nettoyage traite des mesures
  invraisemblables, et déguiser une erreur en absence revient à la
  perdre.
\end{itemize}

Ni l'un ni l'autre ne participent au calcul des seuils.

\section{Quatre méthodes}

\begin{longtable}[]{@{}
  >{\raggedright\arraybackslash}p{(\linewidth - 4\tabcolsep) * \real{0.3333}}
  >{\raggedright\arraybackslash}p{(\linewidth - 4\tabcolsep) * \real{0.3333}}
  >{\raggedright\arraybackslash}p{(\linewidth - 4\tabcolsep) * \real{0.3333}}@{}}
\caption{Méthodes de nettoyage}\label{tab:nettoyage}\\
\toprule\noalign{}
\begin{minipage}[b]{\linewidth}\raggedright
Méthode
\end{minipage} & \begin{minipage}[b]{\linewidth}\raggedright
Seuil
\end{minipage} & \begin{minipage}[b]{\linewidth}\raggedright
À utiliser quand
\end{minipage} \\
\midrule\noalign{}
\endhead
\bottomrule\noalign{}
\endlastfoot
\texttt{RemoveOutbounds(min,\ max)} & Fixe & Bornes fixes: une pluie négative, une température au-dessus de
l'ébullition \\
\texttt{RemovePercentileOutliers(low,\ high)} & Percentiles des relevés
& Elimination des queues de distribution \\
\texttt{RemoveZScoreOutliers(level)} & Moyenne ± level × écart-type &
Les relevés se dispersent symétriquement ; 3 est l'usage \\
\texttt{RemovePeirceOutliers()} & Critère de Peirce & Aucun seuil ne
s'impose \\
\end{longtable}

Les bornes s'expriment en \texttt{float64}, avec \texttt{NaV} pour « pas
de borne de ce côté » :

\begin{verbatim}
ts.RemoveOutbounds(0, nav.NaV)            // rien en dessous de zéro
ts.RemovePercentileOutliers(nav.NaV, 95)  // couper le haut seulement
\end{verbatim}

\textbf{Sur les percentiles :} ils décrivent les relevés dont on
dispose, donc les seuils bougent avec les données et quelque chose est
toujours rejeté. C'est un outil différent des bornes fixes.


\textbf{Sur le critère de Peirce :} il ne prend aucun seuil. Il le
déduit de la taille de l'échantillon, à travers une table publiée par
Benjamin Peirce en 1852~\cite{peirce1852,ross2003}, et décide combien de relevés un échantillon de
cette taille peut légitimement perdre. Il est plus sévère sur les petits
échantillons, rejette au plus neuf relevés quelle que soit la taille, et
--- le test qui compte --- ne rejette rien du tout d'une série saine.
Sur l'exemple classique d'enseignement, il condamne exactement les deux
relevés bas, 90 et 89.

\chapter{Régularisation}\label{chap:regularisation}

\section{Grille}

\begin{verbatim}
hourly, err := ts.Regularize(time.Hour, timeseries.AggMean)
\end{verbatim}

Une série régularisée est une série dont les temps d'horodatage sont espacés régulièrement. La durée entre l'horodatage précédent et l'horodatage courant forme un 
intervalle semi-ouvert à gauche (fermé à droite).

Les fonctions de régularisation créent effectivement ce qui s'appelle une \textit{série temporelle}, en ce sens que chaque mesure est séparée des autres par une durée identifiable.
On va parler de moyenne quotidienne, de prix à la fin du mois, etc.
Les fenêtres ferment à droite et s'alignent sur l'horloge (\S\ref{sec:utc}).

Quand l'application reçoit plusieurs mesures dans la fenêtre demandée, un choix doit être fait sur le traitement des mesures intermédiaires. Une collection d'agrégateurs a été mise en place à cet effet (§\ref{sec:agregateurs}, table~\ref{tab:agregateurs}).

Une fenêtre qui n'a rien reçu devient un trou : une grille régulière a un
point par pas, et un pas sans relevé reste un pas. Rien n'est émis avant
le premier relevé ni après le dernier --- une série ne dit rien de ce
qui précède son début ni de ce qui suit sa fin.

\section{Tolérance}

Un enregistreur censé rendre son relevé à l'heure pile et qui le rend à
10:00:04 n'a pas manqué sa fenêtre. Sans tolérance, il tombe une fenêtre
trop tard :

\begin{longtable}[]{@{}lll@{}}
\caption{Effet d'une tolérance de cinq minutes}\label{tab:tolerance}\\
\toprule\noalign{}
Fenêtre & Sans tolérance & 5 minutes de tolérance \\
\midrule\noalign{}
\endhead
\bottomrule\noalign{}
\endlastfoot
01:00 & 57 & 57 \\
02:00 & \textbf{NaV} & 123 \\
03:00 & \textbf{151,5} --- la moyenne de deux relevés & 180 \\
04:00 & \textbf{NaV} & 241 \\
\end{longtable}

Le signal horaire devient une alternance de trous et de moyennes
fausses. La tolérance doit rester inférieure au pas, sans quoi un relevé
appartiendrait à deux fenêtres.

\section{Agrégateurs}\label{sec:agregateurs}

\begin{longtable}[]{@{}
  >{\raggedright\arraybackslash}p{(\linewidth - 4\tabcolsep) * \real{0.3333}}
  >{\raggedright\arraybackslash}p{(\linewidth - 4\tabcolsep) * \real{0.3333}}
  >{\raggedright\arraybackslash}p{(\linewidth - 4\tabcolsep) * \real{0.3333}}@{}}
\caption{Agrégateurs de fenêtre}\label{tab:agregateurs}\\
\toprule\noalign{}
\begin{minipage}[b]{\linewidth}\raggedright
Agrégateur
\end{minipage} & \begin{minipage}[b]{\linewidth}\raggedright
Renvoie
\end{minipage} & \begin{minipage}[b]{\linewidth}\raggedright
Exemple d'utilisation
\end{minipage} \\
\midrule\noalign{}
\endfirsthead
\caption[]{Agrégateurs de fenêtre \emph{(suite)}}\\
\toprule\noalign{}
\begin{minipage}[b]{\linewidth}\raggedright
Agrégateur
\end{minipage} & \begin{minipage}[b]{\linewidth}\raggedright
Renvoie
\end{minipage} & \begin{minipage}[b]{\linewidth}\raggedright
Exemple d'utilisation
\end{minipage} \\
\midrule\noalign{}
\endhead
\bottomrule\noalign{}
\endlastfoot
\texttt{AggMean} & La moyenne de la fenêtre & Une grandeur physique \\
\texttt{AggMedian} & La médiane & Idem, quand un relevé égaré ne doit
pas tirer le résultat \\
\texttt{AggMin}, \texttt{AggMax} & Les extrêmes & Mesure de composants toxiques (CO2, CO,...) \\
\texttt{AggSum} & La somme & Une grandeur qui s'accumule : pluie,
énergie, événements \\
\texttt{AggFirst}, \texttt{AggLast} & Le relevé du bord, trou compris &
Un signal d'état, un compteur, le prix d'un bien \\
\texttt{AggFirstUsable}, \texttt{AggLastUsable} & Le relevé du bord, en
enjambant les trous & Une grandeur physique \\
\texttt{AggCountUsable} & Combien de relevés sont arrivés & Un rapport
de couverture \\
\texttt{AggSlope} & La pente de la droite ajustée & Une tendance par
pas \\
\texttt{AggIntegral(step)} & L'aire sous la fenêtre & Transformer un
débit en quantité \\
\end{longtable}

Tous suivent la politique NaV : un trou est ignoré, une erreur se
propage. Une fenêtre vide n'atteint jamais l'agrégateur --- elle est
émise en NaV directement.

\texttt{Aggregator("maximum")} en renvoie un par son nom, pour qu'une
recette stockée en base puisse le choisir. Un nom inconnu est une
erreur, jamais un repli silencieux sur la moyenne.

\chapter{Regroupement calendaire}\label{chap:calendrier}

\texttt{DownscaleDaily}, \texttt{DownscaleWeekly},
\texttt{DownscaleMonthly} et \texttt{DownscaleYearly} regroupent par
période du calendrier.

Une durée constante est satisfaisante d'un point de vue statistique, mais parfois surprenantes à lire sur un calendrier. Car dans le calendrier \textit{social} auquel
nous sommes habitués, rien n'est constant: à part les secondes, tous nos repères habituels (jours, mois, années) sont irréguliers. D'où la nécéssité de faire des regroupements calendaires.


\begin{itemize}
\tightlist
\item
  un mois dure de 28 à 31 jours --- février contient environ 10 \% de
  temps en moins que mars ;
\item
  une année dure 365 ou 366 jours ;
\item
  une journée dure 24 heures, sauf les deux où les horloges changent (\textit{seconde intercalaire} - \textit{leap second}) .
\end{itemize}

Deux points mensuels ne pèsent donc pas la même chose. Un agrégateur qui
croît avec la durée de la fenêtre --- \texttt{AggSum},
\texttt{AggCountUsable}, \texttt{AggIntegral} --- produit des nombres
qui diffèrent en partie parce que les périodes diffèrent, et février
ressort plus bas pour la seule raison qu'il est plus court.

\textbf{Régulariser pour calculer, regrouper en dernier, pour montrer.}

\chapter{Compression et restitution}\label{chap:reduction}

\section{Reduce}

Un signal qui tient sa valeur entre deux changements --- une porte, une
consigne, une machine à états --- stocke le même nombre encore et
encore. \texttt{Reduce} garde le premier point, le dernier, et les
changements entre les deux.

Ce qui compte comme changement suit la politique plutôt qu'une
comparaison ordinaire, puisqu'aucun NaN n'est égal à lui-même : deux
trous consécutifs ne sont pas un changement, mais entrer dans un trou ou
en sortir en est un --- et c'est souvent l'information la plus
intéressante qu'une série enregistre.

\texttt{ReduceWithDeadband(band)} fait de même avec une tolérance : un
relevé n'est gardé que s'il diffère \textbf{du dernier relevé gardé} de
plus que la bande --- et non du relevé précédent, si bien qu'une dérive
lente est attrapée dès qu'elle a bougé de plus que la bande au total.
Celle-là perd de l'information, d'une quantité bornée, et le dit dans
son nom.

\section{Expand}

\texttt{Expand(start,\ end,\ step)} reconstruit une grille régulière en
tenant chaque valeur jusqu'au changement suivant. Avant le premier point
connu, il rend NaV : la série ne dit rien de ce que faisait le signal à
ce moment-là, et un trou vaut mieux qu'une valeur inventée.

\section{MarkSilences}

Réduire une série brute comporte un piège : si l'instrument cesse
d'émettre pendant une journée, rien dans la série réduite n'enregistre
le silence --- la dernière valeur semble simplement se maintenir.

\texttt{MarkSilences(maxGap)} insère un NaV partout où deux relevés sont
espacés de plus de \texttt{maxGap}, daté \texttt{maxGap} après le
dernier relevé précédant le silence --- le moment où une alarme
surveillant le capteur se serait déclenchée. Jusque-là, la valeur tenue
est légitime : un capteur qui émet toutes les dix minutes n'est pas
perdu à la onzième. Le silence survit alors à la réduction, sans passer
par une grille.

\chapter{Interpolation}\label{chap:interpolation}

Seule opération qui invente des données : ce qui en sort est une lecture
de ce qui s'est probablement passé, pas un relevé.

\section{Sur le temps, pas sur les rangs}

Un trou est comblé d'après sa place \textbf{dans le temps}. Sur une
série régularisée, ça ne change rien ; sur une série brute, ça change
tout :

\begin{verbatim}
10 mesuré à 00:00, un trou à 00:01, 20 mesuré à 01:00

   par le temps (cette bibliothèque) : 10,17
   par les rangs (l'ancienne)        : 15
\end{verbatim}

Le trou est à une minute de son voisin de gauche et à cinquante-neuf de
celui de droite. Quinze n'est pas une réponse défendable.

\section{Sept méthodes}

\begin{longtable}[]{@{}
  >{\raggedright\arraybackslash}p{(\linewidth - 4\tabcolsep) * \real{0.3333}}
  >{\raggedright\arraybackslash}p{(\linewidth - 4\tabcolsep) * \real{0.3333}}
  >{\raggedright\arraybackslash}p{(\linewidth - 4\tabcolsep) * \real{0.3333}}@{}}
\caption{Méthodes d'interpolation}\label{tab:interpolation}\\
\toprule\noalign{}
\begin{minipage}[b]{\linewidth}\raggedright
Méthode
\end{minipage} & \begin{minipage}[b]{\linewidth}\raggedright
Comble avec
\end{minipage} & \begin{minipage}[b]{\linewidth}\raggedright
Pour
\end{minipage} \\
\midrule\noalign{}
\endhead
\bottomrule\noalign{}
\endlastfoot
\texttt{InterpLinear} & Une droite dans le temps & Une grandeur qui
varie continûment \\
\texttt{InterpNearest} & Le voisin le plus proche & Un signal en marches
; garde une valeur mesurée \\
\texttt{InterpForwardFill} & La dernière valeur connue & Une consigne,
un état, un compteur \\
\texttt{InterpBackwardFill} & La valeur connue suivante & Le début d'une
série \\
\texttt{InterpLogLinear} & Un taux de croissance constant & Une grandeur
qui se compose ; les deux voisins doivent être positifs \\
\texttt{InterpCubicSpline} & Une spline cubique naturelle & La
régularité avant tout \\
\texttt{InterpMonotoneSpline} & Une spline PCHIP~\cite{fritsch1980} & La régularité sans
dépassement \\
\end{longtable}

\textbf{Pourquoi la spline monotone existe.} À travers une marche de 0 à
10, la spline naturelle comble le trou suivant avec \textbf{14,44} ---
une valeur supérieure à tous les relevés de la série, qu'aucun
instrument n'a vue. La spline monotone reste à 10. Sur des données
mesurées, c'est ce qui en fait la plus sûre des deux.

\section{Seulement les trous, et seulement les courts}

Les valeurs cassées sont laissées telles quelles (\S\ref{sec:regles}), et rien
n'extrapole : un trou sans relevé d'un côté reste un trou, à l'exception
des deux reports, qui s'appuient sur un seul voisin par construction.

Aucune méthode ne sait quelle durée de panne peut être franchie. Une
demi-heure de température manquante se traverse ; trois jours, non ---
et un graphique qui cache la panne derrière une courbe lisse ment sans
en avoir l'air :

\begin{verbatim}
filled, err := ts.InterpolateWithin(timeseries.InterpLinear, 2*time.Hour)
\end{verbatim}

La limite se mesure entre les relevés qui encadrent le trou : elle
signifie donc la même chose sur une série brute et sur une grille.

\chapter{Conteneurs}\label{chap:conteneurs}

Un \texttt{TsContainer} garde les variantes d'un même signal :
\texttt{raw}, \texttt{cleaned}, \texttt{hourly}, tout ce qu'une recette
a produit.

Il les conserve \textbf{dans l'ordre où elles ont été déposées}, ce
qu'une map Go ne fait pas : son parcours est aléatoire. La version
précédente imprimait donc les variantes dans un ordre différent à chaque
exécution, et une légende de graphique se remélangeait entre deux
appels. Remplacer une variante lui conserve son rang.

\texttt{Get} distingue un nom jamais déposé d'un nom déposé à
\texttt{nil} : le premier n'a jamais été demandé, le second n'a pas pu
être produit, et l'affichage dit lequel.

\chapter{Sorties}\label{chap:sorties}

\section{Terminal}

\begin{verbatim}
ts.PrettyPrint()       // le tableau des points
ts.PrintStats()        // le résumé, en quatre sections
ts.PrettyPrint(90, 95) // une fenêtre, pour une longue série
\end{verbatim}

Les sentinelles s'affichent par leur nom --- \texttt{NaV}, \texttt{NaN},
\texttt{NaDuration}, un tiret pour un instant inconnu. Un trou affiché «
NaN », ou sous la forme d'une date de l'an 1, est un trou que personne
ne remarque.

Les variantes \texttt{Fprint*} prennent un \texttt{io.Writer}, pour un
journal ou un test.

\section{JSON}\label{sec:json}

\texttt{ToJSON} produit la forme que lit un frontend : des colonnes
plutôt que des objets.

\begin{longtable}[]{@{}
  >{\raggedright\arraybackslash}p{(\linewidth - 2\tabcolsep) * \real{0.5000}}
  >{\raggedright\arraybackslash}p{(\linewidth - 2\tabcolsep) * \real{0.5000}}@{}}
\caption{Champs de l'export JSON}\label{tab:json}\\
\toprule\noalign{}
\begin{minipage}[b]{\linewidth}\raggedright
Clé
\end{minipage} & \begin{minipage}[b]{\linewidth}\raggedright
Contenu
\end{minipage} \\
\midrule\noalign{}
\endhead
\bottomrule\noalign{}
\endlastfoot
\texttt{chron} & Les instants, en RFC 3339 \\
\texttt{meas}, \texttt{dmeas} & Mesures et variations, les non-nombres
en \texttt{null} \\
\texttt{dchron\_ns} & Les intervalles en nanosecondes,
\texttt{NaDuration} en \texttt{null} \\
\texttt{stats} & Le résumé, mêmes noms de champs que
\texttt{BasicStats}, en minuscules \\
\end{longtable}

Tout non-nombre devient \texttt{null}, JSON n'ayant pas de NaN.
\textbf{La distinction entre un trou et une erreur ne survit donc pas à
la conversion} --- ce sont les compteurs \texttt{nbreOfNaV} et
\texttt{nbreOfNaN} du résumé qui la transportent.

\chapter{Performances, mesurées}\label{chap:performances}

Toutes les durées citées dans ce guide --- la table~\ref{tab:performances}
ci-dessous, mais aussi les tables~\ref{tab:extraction}
et~\ref{tab:chargement} --- ont été relevées sur la même machine : un
Apple Mac mini M4 Pro, douze cœurs (huit de performance, quatre
d'efficacité), 48 Gio de mémoire.

Les valeurs absolues n'ont donc pas de portée générale. Elles dépendent
de la machine, de la version du compilateur, et de ce que celle-ci
faisait par ailleurs. Ce qui porte le raisonnement, ce sont les rapports
entre elles : cinquante contre un entre une médiane et une moyenne, un
ordre de grandeur entre \texttt{Add} et \texttt{AddBatchData}.

\begin{longtable}[]{@{}ll@{}}
\caption{Coûts mesurés, par opération}\label{tab:performances}\\
\toprule\noalign{}
Opération & Coût \\
\midrule\noalign{}
\endhead
\bottomrule\noalign{}
\endlastfoot
Un point en mémoire & 32 octets \\
Moyenne, min, max, somme & \textasciitilde1,2 ns par point, sans
allocation \\
Écart-type & \textasciitilde2 ns par point, sans allocation \\
Médiane, percentile & \textasciitilde50 ns par point --- ils trient une
copie \\
Un \texttt{Stats()} complet & \textasciitilde60 ns par point \\
\texttt{AddBatchData}, dans l'ordre & \textasciitilde0,03 µs par
point \\
\texttt{AddBatchData}, en désordre & \textasciitilde0,5 µs par point \\
\texttt{Add}, un par un, en désordre & quadratique --- à éviter \\
\end{longtable}

La croissance est linéaire jusqu'à dix millions de points pour tout ce
qui n'est pas médiane, celles-ci portant leur \texttt{n\ log\ n}.

\chapter{Migration depuis la version précédente}\label{chap:migration}

La bibliothèque a été réécrite pour la \texttt{v0.2}. Ce qui a changé,
et pourquoi :

\begin{longtable}[]{@{}
  >{\raggedright\arraybackslash}p{(\linewidth - 4\tabcolsep) * \real{0.3333}}
  >{\raggedright\arraybackslash}p{(\linewidth - 4\tabcolsep) * \real{0.3333}}
  >{\raggedright\arraybackslash}p{(\linewidth - 4\tabcolsep) * \real{0.3333}}@{}}
\caption{Changements apportés par la v0.2}\label{tab:migration}\\
\toprule\noalign{}
\begin{minipage}[b]{\linewidth}\raggedright
Avant
\end{minipage} & \begin{minipage}[b]{\linewidth}\raggedright
Maintenant
\end{minipage} & \begin{minipage}[b]{\linewidth}\raggedright
Pourquoi
\end{minipage} \\
\midrule\noalign{}
\endhead
\bottomrule\noalign{}
\endlastfoot
Un point rejeté devenait \texttt{NaN} & Il devient \texttt{NaV} & Une
valeur aberrante ne détruit plus un mois de statistiques \\
Statistiques stockées dans la série & \texttt{Stats()} calcule à la
demande & Des statistiques stockées se périmaient en silence \\
Un drapeau \texttt{deltasValid} & Un invariant & Rendre le mauvais état
impossible vaut mieux que le signaler \\
\texttt{DataSeries} public & Privé, avec des accesseurs & Ce drapeau
n'existait que parce que le tableau était public \\
Tolérance acceptée, jamais appliquée & Appliquée & Elle était
silencieusement ignorée \\
Interpolation par rangs & Par le temps & Le rang est faux sur toute
série irrégulière \\
Interpolation modifiant en place & Rend une nouvelle série & Cohérent
avec le nettoyage et la régularisation \\
\texttt{MemId\ uint64} & \texttt{ID\ string} & Il n'était jamais rempli
; une chaîne accepte une clé de base ou un UUID \\
\texttt{Chstd} & Supprimé & La dispersion des horodatages ne dit rien ;
\texttt{DChstd} dit l'utile \\
--- & \texttt{ChFirstUsable} & Là où la fenêtre s'ouvre n'est pas là où
les données commencent \\
--- & \texttt{NbreOfNaV} & Distinguer les trous des erreurs \\
--- & \texttt{InterpolateWithin}, \texttt{MarkSilences},
\texttt{ReduceWithDeadband} & Nouveaux \\
\end{longtable}

Pour le contrat JSON : \texttt{id} est désormais une chaîne,
\texttt{chstd} disparaît, et \texttt{nbreOfNaV}, \texttt{chFirstUsable}
et \texttt{valAtFirstUsable} s'ajoutent.

\chapter{Traçabilité des décisions}\label{chap:raisonnement}

Chaque décision ci-dessus est soutenue par un test qui l'énonce en
toutes lettres et par un message de commit qui consigne pourquoi elle a
été prise. Quand un comportement surprend, \texttt{git\ log} et
\texttt{git\ blame} sur le fichier l'expliquent mieux que le code.

% Les annexes : les exemples d'utilisation et la liste des signatures.
% Leur code est extrait des paquets par docs/gendoc.py.
% Les annexes du guide.
%
% Le code qu'elles présentent n'est pas saisi ici : docs/gendoc.py
% l'extrait des paquets à chaque composition, dans docs/generated/.
% Ce fichier ne contient que la prose et l'ordre de présentation.
%
% Ajouter un exemple demande donc deux gestes : la fonction Example
% dans example_test.go, et le \input correspondant ci-dessous. Si le
% second manque, l'exemple n'est pas dans le guide ; s'il désigne un
% exemple disparu, la composition s'arrête et le dit.

\appendix

\chapter{Exemples d'utilisation}\label{chap:exemples}

Les exemples qui suivent ne sont pas des extraits recopiés. Ce sont des
fonctions \texttt{Example} des deux paquets : \texttt{go\ test} les
compile, les exécute, et compare ce qu'elles écrivent au commentaire
\texttt{//\ Output} qui les termine. Un exemple qui ne correspond plus
au code fait échouer les tests.

C'est la seule forme de documentation qui ne puisse pas vieillir sans
qu'on le sache, et la raison pour laquelle ce qui suit est extrait de
\texttt{example\_test.go} au moment de composer le guide plutôt que
saisi une fois pour toutes. Les mêmes exemples se déplient sur
pkg.go.dev, sous la fonction qu'ils illustrent.

Ils sont écrits dans le paquet de test externe : ils montrent donc
l'import et le préfixe qu'un appelant écrit réellement, et non les
raccourcis que seul le paquet lui-même s'autorise. Les commentaires
sont en anglais, comme le code.

Deux fonctions auxiliaires, \texttt{at} et \texttt{show}, abrègent les
exemples de \texttt{timeseries} : la première fabrique un relevé pris
\texttt{n} minutes après 10:00 UTC, la seconde affiche une série à
raison d'une ligne par relevé.

\section{\texttt{notavalue}}

\subsection{Les trois règles sur un même calcul}

Un relevé manquant est ignoré, un calcul rompu se propage, et un
ensemble où il n'y a rien à calculer ne donne pas zéro mais rien.
Chapitre~\ref{chap:nav}.

\input{generated/exemple-notavalue-Example}

\subsection{L'asymétrie de l'addition}

L'addition traite le manque comme neutre, les trois autres opérations
non. Une pluie mensuelle dont un jour n'a pas été relevé reste la pluie
tombée ; une concentration manquante multipliée par un débit connu ne
donne rien.

\input{generated/exemple-notavalue-ExampleAdd}

\subsection{Choisir la première source disponible}

Un capteur principal, son secours, puis une valeur modélisée.
\texttt{Coalesce} passe sur ce qui manque, mais pas sur une erreur : un
calcul rompu en amont ne se répare pas en passant au plan B.

\input{generated/exemple-notavalue-ExampleCoalesce}

\subsection{Ce que vaut un résultat}

Une moyenne ne dit pas sur combien de relevés elle repose. Les
compteurs le disent, et distinguent les relevés absents des calculs
rompus --- ces derniers sont à chercher en amont.

\input{generated/exemple-notavalue-ExampleCountUsable}

\subsection{Garder la distinction à l'écran}

Les verbes de la bibliothèque standard voient les deux sentinelles
comme des NaN et écrivent « NaN » pour l'une comme pour l'autre. Toute
la distinction disparaît dans une ligne de journal ou un message
d'échec de test ; \texttt{Format} la conserve.
Chapitre~\ref{chap:sorties}.

\input{generated/exemple-notavalue-ExampleFormat}

\subsection{Une médiane qui n'invente pas de valeur}

Sur un nombre pair de relevés, cette médiane rend la plus petite des
deux valeurs centrales plutôt que leur moyenne. Sur une vanne codée 0
et 1, la définition classique renverrait 0,5 --- un état que
l'équipement n'a jamais occupé.

\input{generated/exemple-notavalue-ExampleMedian}

\section{\texttt{timeseries}}

\subsection{Le parcours complet}

Des relevés bruts à un résumé : les valeurs invraisemblables écartées,
une grille régulière, les trous courts comblés. Les relevés sont remis
dans n'importe quel ordre, ce qui ne coûte rien --- chapitre
\ref{chap:serie}.

\input{generated/exemple-timeseries-Example}

\subsection{Un trou et une erreur ne sont pas le même accident}

La première série a un relevé manquant, et garde une moyenne. La
seconde porte le résultat d'un calcul rompu, et n'en a plus. La même
asymétrie se lit sur la variation d'un point au suivant : un relevé de
20 après un trou n'est pas une hausse de 20.
Chapitre~\ref{chap:resume-stats}.

\input{generated/exemple-timeseries-Example-gapVersusError}

\subsection{Nettoyer, et garder ce qu'on a rejeté}

Le nettoyage rend deux séries : celle qui est nettoyée et celle des
relevés écartés, pour que la décision puisse être revue. Dans la série
nettoyée, l'instant rejeté est toujours là --- sa mesure est devenue
manquante, elle n'a pas disparu. Chapitre~\ref{chap:nettoyage}.

\input{generated/exemple-timeseries-ExampleTimeSeries-RemoveOutbounds}

\subsection{Un enregistreur qui dérive}

Un enregistreur censé rapporter à l'heure juste rapporte quelques
secondes trop tard. Sans tolérance, chacun de ces relevés tombe dans la
fenêtre suivante, ce qui laisse une fenêtre vide et met deux relevés
dans la suivante. Chapitre~\ref{chap:regularisation}.

\input{generated/exemple-timeseries-ExampleTimeSeries-RegularizeWithTolerance}

\subsection{Combler un silence court, pas une panne}

La limite se mesure entre les deux relevés qui encadrent le trou, et
non entre les pas de la grille : elle veut donc dire la même chose sur
une série régulière et sur une série brute.
Chapitre~\ref{chap:interpolation}.

\input{generated/exemple-timeseries-ExampleTimeSeries-InterpolateWithin}

\subsection{Comprimer un signal d'état, puis le restituer}

Un signal qui passe son temps à ne pas bouger. \texttt{Reduce} garde
les instants où quelque chose s'est produit, \texttt{Expand} remet la
grille. Chapitre~\ref{chap:reduction}.

\input{generated/exemple-timeseries-ExampleTimeSeries-Reduce}

\subsection{Parcourir un conteneur sans allouer}

Un conteneur tient les séries qui vont ensemble --- les relevés bruts,
la version nettoyée, chaque variante produite en chemin.
\texttt{MeasTo} est le compagnon de la boucle intérieure : on lui
confie une tranche de longueur nulle, il la remplit au lieu d'en
allouer une neuve, et un passage sur cent séries n'alloue qu'une fois.
Chapitre~\ref{chap:conteneurs}.

\input{generated/exemple-timeseries-ExampleTsContainer}

\chapter{Référence des interfaces}\label{chap:api}

Les listes qui suivent donnent, pour chaque paquet, ce qu'il exporte :
constantes, variables, fonctions, types avec leurs champs et leurs
méthodes. Elles sont extraites de \texttt{go\ doc\ -all} et ne peuvent
donc décrire ni une fonction qui n'existe plus, ni oublier celle qui
vient d'être ajoutée.

Ce sont des signatures, sans les commentaires qui les accompagnent. Le
godoc complet est sur pkg.go.dev, pour
\href{https://pkg.go.dev/github.com/fflamingodev/notavalue}{\texttt{notavalue}}~\cite{pkgnotavalue}
et pour
\href{https://pkg.go.dev/usefulrisk.com/timeseries}{\texttt{timeseries}}~\cite{pkgtimeseries},
où il est navigable et à jour de la version publiée. Le reproduire ici
doublerait ce document pour répéter en anglais ce que les chapitres
précédents expliquent en français.

L'ordre est celui de \texttt{go\ doc} : alphabétique à l'intérieur de
chaque rubrique, les constructeurs et les méthodes d'un type venant
après sa déclaration.

\input{generated/api-notavalue}

\input{generated/api-timeseries}


\printbibliography[heading=bibintoc, title={Bibliographie}]

\end{document}
